Il lavoro ha realizzato \emph{Copernicus Data Workspace}: un assistente per
consultare fonti geospaziali tramite richieste in linguaggio naturale, esposto
in una dashboard web autenticata con Keycloak. Il contributo è progettuale e
implementativo: responsabilità specializzate e percorsi controllati. Sul piano
progettuale sono stati ottenuti, in particolare:

\begin{itemize}
    \item un workflow LangGraph con confine di dominio, parsing strutturato,
    percorso Discovery su NebulaGraph, percorso operativo multi-sorgente
    (web, GDACS, STAC live, CEMS) e percorsi combinati;
    \item integrazioni MCP per GDACS e CEMS, client STAC diretto per le
    acquisizioni, harvester e catalogo metadati Marine/STAC/Drought;
    \item memoria conversazionale Redis, artefatti dashboard verificabili
    (valori calcolati dal codice sui record recuperati), ripresa delle
    sessioni e interfaccia bilingue IT/EN;
    \item stack di osservabilità (Prometheus, Grafana, OpenTelemetry/Phoenix)
    per tempi, errori, nodi e dipendenze;
    \item un protocollo di valutazione del workflow
    (versione~6.0) con controlli deterministici, rubriche ordinali e stima
    dell'incertezza.
\end{itemize}

MCP delimita l'integrazione di GDACS e CEMS; CDSE STAC live resta un percorso
diretto. Discovery e STAC live svolgono compiti distinti: il primo sceglie
risorse dal catalogo a grafo; il secondo cerca scene sul catalogo remoto.

\textbf{Risultati della valutazione.}
Sul dataset descritto da 105 casi, il riepilogo della run di riferimento del
protocollo~6.0 (Judge \texttt{gemma3:27b}, temperatura~0) riporta
sul grafo 31/49 (63,3\% primario) e 31/51 (60,8\% rigoroso), e sulle
acquisizioni 54/54 (100\%). I conteggi e il ricalcolo degli intervalli sono documentati nel rapporto di
riscontro tecnico del 5 novembre 2026 \cite{riscontroTecnico2026}.
Il Capitolo~\ref{cap:7} ne esplicita criteri e denominatori.
Il benchmark copre questi due percorsi; GDACS, CEMS, follow-up e ripresa
sessione sono oggetto di altre verifiche.

Una seconda run con seed diverso ottiene valori vicini (61,2\% e 98,1\%).
Il recupero sul grafo è buono, ma il riferimento è raramente al primo posto;
sulle acquisizioni il benchmark è ormai troppo facile per distinguere
miglioramenti.

\textbf{Limiti e sviluppi.}
Le priorità sono il ranking della Discovery, un benchmark delle acquisizioni
più difficile, la protezione delle API del backend con verifica del token e
del proprietario degli artefatti, una baseline controllata sul contributo
multiagente, la validazione umana del Judge e la verifica nel browser di
risposta e artefatto. Sul piano operativo vanno ampliati
telemetria, freschezza del catalogo e recupero dai guasti. L'obiettivo è
mantenere ogni affermazione proporzionata all'evidenza disponibile.
